\documentclass{amsart}
\usepackage{amsmath,amssymb,amsthm,hyperref}
\usepackage{algorithm}
\usepackage{algpseudocode}

\newtheorem{theorem}{Theorem}
\newtheorem{lemma}[theorem]{Lemma}
\newtheorem{proposition}[theorem]{Proposition}
\newtheorem{corollary}[theorem]{Corollary}
\theoremstyle{definition}
\newtheorem{definition}[theorem]{Definition}
\newtheorem{remark}[theorem]{Remark}

\begin{document}

\title{On the competitive ratio of the Drop-The-Larger-Hopeless policy}
\maketitle

\section{Formalization of the model}

We make the informal statement precise.

\begin{definition}[Instance]
An \emph{instance} is a finite multiset
$\omega=\{(a_j,s_j,d_j)\}_{j=1}^{n}$ of jobs, where $a_j\in\mathbb{R}_{\ge 0}$
is the arrival (release) time, $s_j>0$ the size, and $d_j\ge a_j+s_j$ the
deadline of job $j$. A single server processes work at unit rate and may
preempt at no cost.
\end{definition}

\begin{definition}[Feasibility]\label{def:feas}
Fix a time $t$ and a finite set $Q$ of jobs, each job $j\in Q$ carrying a
\emph{remaining size} $r_j>0$ and a deadline $d_j\ge t$. We say $Q$ is
\emph{feasible at time $t$} if there is a unit-rate single-server schedule of
$Q$ on $[t,\infty)$ completing every $j\in Q$ by $d_j$. Equivalently, by the
exchange (earliest-deadline-first) argument, $Q$ is feasible at $t$ iff, listing
$Q=\{j_1,\dots,j_m\}$ with $d_{j_1}\le\cdots\le d_{j_m}$,
\begin{equation}\label{eq:edf}
t+\sum_{i=1}^{\ell} r_{j_i}\ \le\ d_{j_\ell}\qquad\text{for every }\ell=1,\dots,m.
\end{equation}
A job is \emph{hopeless} at $t$ if it cannot be completed by its deadline under
any continuation, i.e.\ already $t+r_j>d_j$.
\end{definition}

Throughout, the server is allowed to realize \emph{any} feasible schedule of its
currently admitted set; ``feasible'' always refers to
Definition~\ref{def:feas}. (The phrase ``serves in FCFS order'' in the prompt is
only the default tie-break used by DTLH to decide \emph{which jobs are present};
the admission test and the completion guarantee use the feasibility relation
\eqref{eq:edf}, as the prompt states: a set is feasible if there \emph{exists} a
schedule meeting all deadlines.)

\begin{definition}[The DTLH policy]\label{def:dtlh}
DTLH maintains an \emph{admitted set} $A$ that is feasible at all times; an
admitted job is guaranteed completion, a non-admitted job is \emph{dropped} and
counts as one miss. Decisions are made online at arrival epochs. When job $j$
arrives at time $t=a_j$, let $A^-$ be the set of admitted, not-yet-completed
jobs at $t$ (with their remaining sizes). DTLH acts as in
Algorithm~\ref{alg:dtlh}.
\end{definition}

\begin{remark}[The size used in the eviction test]\label{rem:sizeconv}
The prompt says DTLH ``compares the new job's size with the \emph{sizes} of
currently queued jobs,'' where size $=$ service requirement $s_j$. We therefore
take the eviction comparison in line~4 to use the jobs' \emph{original} sizes
$s_k,s_j$, \emph{not} their remaining work $r_k$. This choice is mathematically
material: the same instance can have different DTLH miss counts under the two
readings (Proposition~\ref{prop:five} achieves $5$ misses under the
original-size reading and only $3$ under a remaining-work reading). We adopt the
original-size reading because (i) it is the literal meaning of ``size $=$ service
requirement'' in the prompt, and (ii) ``size'' is a fixed attribute of a job,
known at admission, whereas remaining work is schedule-dependent. The
\emph{feasibility} test in line~1 and line~4, by contrast, necessarily uses
remaining work $r_k$, since it concerns whether the surviving set can still be
completed.
\end{remark}

\begin{algorithm}[h]
\caption{DTLH admission decision at the arrival of job $j$ at time $t$}
\label{alg:dtlh}
\begin{algorithmic}[1]
\Require current admitted, not-yet-completed set $A^-$ (remaining sizes $r_k$, original sizes $s_k$), arriving job $j=(s_j,d_j)$, time $t$
\Ensure updated admitted set and the miss count increment
\If{$A^-\cup\{j\}$ is feasible at $t$ (test \eqref{eq:edf} on remaining sizes)}
  \State admit $j$: \ $A^-\gets A^-\cup\{j\}$
\Else
  \State $L\gets\{k\in A^- : s_k> s_j \text{ and } (A^-\setminus\{k\})\cup\{j\}\text{ is feasible at }t\}$
  \If{$L\neq\varnothing$}
    \State pick any $k^\star\in L$; drop $k^\star$, admit $j$: \ $A^-\gets(A^-\setminus\{k^\star\})\cup\{j\}$; one miss ($k^\star$)
  \Else
    \State drop $j$; one miss ($j$)
  \EndIf
\EndIf
\State \Return $A^-$
\end{algorithmic}
\end{algorithm}

\begin{remark}[The policy is underspecified; we measure its worst case]
\label{rem:tiebreak}
Algorithm~\ref{alg:dtlh} contains two genuine sources of nondeterminism: the
order in which simultaneous arrivals (same $a_j$) are processed, and the choice
of victim $k^\star\in L$ in line~5. The prompt fixes neither (it says ``drops
\emph{such} a larger queued job''). We therefore regard DTLH as a \emph{family}
of policies indexed by these tie-breaks, and we define
\[
\mathrm{Miss}_{\mathrm{DTLH}}(\omega)
=\max\{\text{number of jobs dropped over all legal DTLH tie-break choices on }\omega\},
\]
the \emph{worst-case} behaviour of the rule. This is the correct quantity for a
\emph{lower} bound on the competitive ratio of the policy class: a single legal
execution achieving many misses already witnesses that the rule, as stated, is
not better than the ratio so obtained. (A specific deterministic tie-break can do
strictly better on individual instances; indeed, as we note below, on every
lower-bound instance of this paper there \emph{is} a legal tie-break that does
better. Our lower-bound instances are therefore presented together with the
explicit legal execution that realizes the stated worst-case miss count, so the
bounds hold for the worst-case reading and a fortiori for the supremum over the
policy class.)
\end{remark}

\begin{definition}[Benchmarks]
For an instance $\omega$, $\mathrm{Miss}_{\mathrm{DTLH}}(\omega)$ is as in
Remark~\ref{rem:tiebreak}. Let
\begin{equation}\label{eq:opt}
\mathrm{Miss}_{\mathrm{OPT}}(\omega)
=\min\Big\{\,|\omega\setminus S|\ :\ S\subseteq\omega \text{ is feasible respecting release times}\,\Big\},
\end{equation}
the minimum number of jobs that any policy must miss; $S$ feasible respecting
release times means: the preemptive single-machine EDF schedule of $S$ (run, at
each instant, an available unfinished job of earliest deadline) completes every
job of $S$ by its deadline. By the optimality of EDF for the feasibility of a
fixed job set with release times, $S$ is schedulable iff its EDF schedule meets
all deadlines, so \eqref{eq:opt} is well defined. This
$\mathrm{Miss}_{\mathrm{OPT}}$ is the clairvoyant offline optimum and is a lower
bound on the misses of every (online or offline) admissible policy; using it in
the denominator therefore \emph{upper-bounds} the achievable competitive ratio,
which is the strongest way to test optimality of DTLH.
\end{definition}

The headline quantity is
\(
R \;=\; \sup_{\omega:\ \mathrm{Miss}_{\mathrm{OPT}}(\omega)>0}
\frac{\mathrm{Miss}_{\mathrm{DTLH}}(\omega)}{\mathrm{Miss}_{\mathrm{OPT}}(\omega)} .
\)

\section{Two structural facts}

\begin{lemma}[No needless drops]\label{lem:nowaste}
If $\omega$ is feasible respecting release times, i.e.\
$\mathrm{Miss}_{\mathrm{OPT}}(\omega)=0$, then \emph{every} legal DTLH execution
drops nothing: $\mathrm{Miss}_{\mathrm{DTLH}}(\omega)=0$.
\end{lemma}

\begin{proof}
We show by induction over arrival epochs that DTLH's admitted set equals the set
of all already-arrived, not-yet-completed jobs (no job has been dropped), and is
feasible at the current time. Initially $A=\varnothing$. Suppose at the arrival
of job $j$ at $t=a_j$ the admitted, not-yet-completed set $A^-$ consists of
exactly the arrived-but-uncompleted jobs. Because $\omega$ is feasible
respecting release times, the EDF schedule of $\omega$ completes all of $\omega$;
restricting that schedule to $[t,\infty)$ yields a feasible completion at time
$t$ of all jobs that have arrived by $t$ and are not yet finished \emph{together
with} $j$ (every such job is in $\omega$, has deadline $\ge t$, and the global
EDF schedule meets its deadline). Hence $A^-\cup\{j\}$ is feasible at $t$, so
DTLH executes line~2 and admits $j$ without any drop, independently of any
tie-break. The invariant is preserved. By induction no drop ever occurs.
\end{proof}

Lemma~\ref{lem:nowaste} shows the ratio is only nontrivial when
$\mathrm{Miss}_{\mathrm{OPT}}(\omega)\ge 1$; it also rules out the degenerate
``divide by zero'' regime.

\begin{lemma}[Single common window: the \emph{evict-largest} rule is optimal]
\label{lem:onewindow}
Suppose all jobs share one release time $a$ and one deadline $d$, so feasibility
of a set $T$ is the single constraint $a+\sum_{j\in T}s_j\le d$, i.e.\ total size
at most the capacity $C:=d-a$. Consider the deterministic DTLH variant
$\mathrm{DTLH}_{\max}$ that, whenever a swap is performed, always evicts a victim
of \emph{largest} original size among the eligible set $L$ (line~4), processing
simultaneous arrivals in any fixed order. Then
$\mathrm{DTLH}_{\max}$ keeps a maximum-cardinality feasible subset, so its miss
count equals $\mathrm{Miss}_{\mathrm{OPT}}(\omega)$.
\end{lemma}

\begin{proof}
Order the jobs by the order in which DTLH processes them, $j_1,j_2,\dots,j_n$
(arrivals are simultaneous, so this is the chosen processing order). Because all
jobs share the window $[a,d]$, every admitted job has deadline $d$ and the
feasibility test \eqref{eq:edf} reduces to ``total admitted size $\le C$,'' while
the eviction comparison ``$s_k>s_{j_i}$'' on original sizes is exactly ``a
strictly larger admitted job exists'' (remaining sizes are irrelevant to a
single-constraint window). Write $A_i$ for the multiset of admitted sizes after
$j_1,\dots,j_i$ have been processed, and let $S_i=\{s_{j_1},\dots,s_{j_i}\}$ be
the multiset of sizes seen so far.

\emph{Invariant.} For every $i$, $A_i$ is a feasible subset of $S_i$ of
\emph{maximum cardinality}, and among all maximum-cardinality feasible subsets of
$S_i$ it has \emph{minimum total size}.

The invariant holds for $i=0$ ($A_0=\varnothing$). Assume it holds for $i-1$ and
process $j_i$ (size $s:=s_{j_i}$). Let $k:=|A_{i-1}|$ and $\sigma:=\sum A_{i-1}$;
by hypothesis $k$ is the maximum feasible cardinality on $S_{i-1}$ and $\sigma$ is
the least total size achieving it.

\emph{Case 1: $\sigma+s\le C$.} DTLH admits $j_i$, so $A_i=A_{i-1}\cup\{s\}$ has
$k+1$ jobs. No feasible subset of $S_i$ can have more than $k+1$ jobs: deleting
one job from such a subset would give a feasible $(\ge k+1)$-subset of either
$S_{i-1}$ (if the deleted job is $j_i$) contradicting maximality $k$, or, if it
contains $j_i$, replacing $j_i$ by nothing still leaves $\ge k$ jobs of
$S_{i-1}$. Hence $k+1$ is maximum. Minimality of the total size: the minimum total
size of a $(k+1)$-subset of $S_i$ is achieved by taking the $k+1$ smallest
elements of $S_i$ whose sum is $\le C$; since the $k$ smallest of $S_{i-1}$ have
sum $\sigma$ (by the inductive minimality) and $s$ is now appended, the running
``smallest-first'' optimum is $A_{i-1}\cup\{s\}$ unless some element of $A_{i-1}$
exceeds $s$ \emph{and} swapping reduces the sum while keeping $k+1$ feasible---but
that cannot happen here because we did \emph{not} drop any job, and the invariant
on $A_{i-1}$ already made it the $k$ smallest; appending $s$ keeps it the $k+1$
smallest (a smallest-first greedy is order-independent). Thus $A_i$ is a
min-sum max-cardinality set.

\emph{Case 2: $\sigma+s>C$.} Admitting outright is infeasible. If some admitted
$k^\star$ has $s_{k^\star}>s$ and $\sigma-s_{k^\star}+s\le C$, then $L\neq
\varnothing$ and $\mathrm{DTLH}_{\max}$ evicts the \emph{largest} admitted size
$M:=\max\{x\in A_{i-1}: x>s,\ \sigma-x+s\le C\}$, giving
$A_i=(A_{i-1}\setminus\{M\})\cup\{s\}$, again $|A_i|=k$. We claim $A_i$ is
min-sum among $k$-subsets of $S_i$ and that no $(k+1)$-subset of $S_i$ is
feasible.

\emph{No bigger cardinality.} If $S_i$ had a feasible $(k+1)$-subset $T$, then
$T\setminus\{j_i\}$ (or $T$ itself, if $j_i\notin T$) is a feasible
$\ge k$-subset of $S_{i-1}$ of size $\le k$ jobs $+$ possibly $j_i$; in either
case $T$ contains at most $k$ elements of $S_{i-1}$, so $j_i\in T$ and
$T\setminus\{j_i\}$ is a feasible $k$-subset of $S_{i-1}$ with total
$\le C-s$. By inductive minimality the least total of a $k$-subset of $S_{i-1}$ is
$\sigma$, so $\sigma\le \sum(T\setminus\{j_i\})\le C-s$, i.e.\ $\sigma+s\le C$,
contradicting Case~2's hypothesis $\sigma+s>C$. Hence $k$ is still maximum.

\emph{Min sum.} The minimum total of a $k$-subset of $S_i$ equals the sum of the
$k$ smallest elements of $S_i$ whose total is $\le C$. Since $A_{i-1}$ is the $k$
smallest of $S_{i-1}$ (inductive minimality), inserting $s$ into the sorted list
and dropping the current largest kept element exactly realizes the new
$k$-smallest-with-sum-$\le C$ set, which is precisely $(A_{i-1}\setminus\{M\})
\cup\{s\}$ when $M$ is the largest kept element exceeding $s$ (and equals
$A_{i-1}$ when $s$ is not among the $k$ smallest, i.e.\ when $L=\varnothing$,
handled next). Thus $A_i$ is min-sum.

If instead $L=\varnothing$ (no admitted job is both larger than $s$ and a
feasibility-restoring victim), then every admitted job is $\le s$ or its removal
does not restore feasibility; since all share the window, ``removal restores
feasibility'' is automatic once $x>s$ (because $\sigma-x+s<\sigma\le C$ would
follow from $x>s$), so $L=\varnothing$ means no admitted job exceeds $s$, i.e.\
$s$ is $\ge$ every kept element; then $s$ is not among the $k$ smallest of $S_i$,
the smallest-first optimum is unchanged, and DTLH correctly drops $j_i$, keeping
$A_i=A_{i-1}$ min-sum and max-cardinality. The invariant is preserved.

By induction $A_n$ is a maximum-cardinality feasible subset of all $n$ jobs, so
$\mathrm{DTLH}_{\max}$ misses $n-|A_n|=\mathrm{Miss}_{\mathrm{OPT}}(\omega)$.
\end{proof}

\begin{remark}[The worst-case tie-break is \emph{not} optimal, already in one
window]\label{rem:onewindowbad}
Lemma~\ref{lem:onewindow} is about the \emph{evict-largest} rule. For the
worst-case quantity $\mathrm{Miss}_{\mathrm{DTLH}}$ of
Remark~\ref{rem:tiebreak} the conclusion \emph{fails}, even in a single common
window: choosing a non-largest eligible victim can cascade into extra drops.
Concretely take capacity $C=5$ (all jobs $a=0$, $d=5$) and the four sizes
$3,2,1,2$ processed in that order. DTLH admits $3$ then $2$ (admitted sizes
$\{3,2\}$, total $5=C$). The size-$1$ arrival overflows ($5+1>5$); both admitted
jobs are larger than $1$ and removable, and the worst-case tie-break evicts the
\emph{smaller} eligible job $2$, giving admitted $\{3,1\}$ (total $4$, miss \#1).
The final size-$2$ arrival overflows ($4+2>5$); the admitted $3>2$ is removable,
so it is evicted, giving $\{1,2\}$ (total $3$, miss \#2). Thus the worst-case run
drops $2$ jobs and keeps only $\{1,2\}$. But the offline optimum keeps the three
jobs $\{1,2,2\}$ (total $5\le C$) and drops only the single size-$3$ job, so
$\mathrm{Miss}_{\mathrm{OPT}}=1$ while
$\mathrm{Miss}_{\mathrm{DTLH}}=2$ here. (Both values are confirmed by exhaustive
enumeration of all legal DTLH tie-breaks.) Hence even on a single window the
\emph{rule class} is at least $2$-competitive, and Lemma~\ref{lem:onewindow}
\emph{cannot} be upgraded to ``every tie-break.'' This already shows that the
constant-factor loss of DTLH is caused by the swap rule's victim choice, not by
multi-deadline interactions; the multi-window gadgets of
Sections~3--\ref{sec:five} amplify the same phenomenon.
\end{remark}

\section{Lower bound: $R\ge 2$}

We exhibit an explicit, verifiable family. Throughout this and the next section,
``$\mathrm{Miss}_{\mathrm{DTLH}}$'' is the worst-case-over-tie-breaks quantity of
Remark~\ref{rem:tiebreak}.

\begin{definition}[Gadget]\label{def:gadget}
The \emph{gadget} $G$ is the three-job instance
\[
G=\{\,B=(a{=}0,\ s{=}3,\ d{=}3),\quad
      W=(a{=}0,\ s{=}3,\ d{=}5),\quad
      H=(a{=}1,\ s{=}1,\ d{=}2)\,\}.
\]
For $c\in\mathbb{Z}_{\ge 0}$ let $G^{(c)}$ be $G$ with every arrival and every
deadline shifted by $6c$, and for $k\ge 1$ set
$\Omega_k=\bigcup_{c=0}^{k-1}G^{(c)}$.
\end{definition}

The copies $G^{(c)}$ occupy pairwise-disjoint time spans
$[6c,\,6c+5]$, so they do not interact; each is its own busy period.

\begin{proposition}\label{prop:gadget}
$\mathrm{Miss}_{\mathrm{DTLH}}(G)=2$ and $\mathrm{Miss}_{\mathrm{OPT}}(G)=1$.
\end{proposition}

\begin{proof}
\emph{Offline optimum.} The subset $\{W,H\}$ is feasible respecting release
times: run $H$ on $[1,2]$ (it arrives at $1$, due $2$) and $W$ on
$[0,1]\cup[2,5]$ (size $3$, arrives at $0$, due $5$); both deadlines are met.
Hence $\mathrm{Miss}_{\mathrm{OPT}}(G)\le 1$. No size-$3$ feasible subset of $G$
exists: keeping all three needs $B$ ($s=3,d=3$) finished by $3$, forcing the
server to spend all of $[0,3]$ on $B$, after which neither $W$ (still $3$ units,
only $[3,5]$ left, $2<3$) nor $H$ (due $2$, already past) can be completed. So
the whole set is infeasible and $\mathrm{Miss}_{\mathrm{OPT}}(G)=1$.

\emph{DTLH (worst-case execution).} We exhibit a legal execution dropping $2$.
Process the two simultaneous $t=0$ arrivals in the order $B$ then $W$.
$A^-=\{B\}$ is feasible ($0+3\le 3$), so $B$ is admitted. Next $W$ arrives at
$t=0$: $\{B,W\}$ has total size $6$; the EDF test \eqref{eq:edf} gives $0+3\le 3$
(for $B$, $d=3$) but $0+3+3=6>5$ (for $W$, $d=5$): infeasible. DTLH looks for
$k\in\{B\}$ with $s_k>s_W=3$; none ($s_B=3\not>3$), so $W$ is dropped (miss
\#1). At $t=1$, $B$ has received one unit of service (remaining $2$, due $3$),
and $H=(s{=}1,d{=}2)$ arrives. The set $\{B,H\}$ at $t=1$ is infeasible: EDF
order is $H$ ($d=2$) then $B$ ($d=3$), giving $1+1=2\le2$ for $H$ but
$1+1+2=4>3$ for $B$. Now $L=\{k\in\{B\}: s_k>s_H=1$ and removing $k$ restores
feasibility$\}$: removing $B$ leaves $\{H\}$, feasible at $t=1$ ($1+1\le 2$), and
$s_B=3>1$, so $B\in L$. DTLH drops $B$ (miss \#2) and admits $H$. This legal run
drops both $B$ and $W$, so $\mathrm{Miss}_{\mathrm{DTLH}}(G)\ge 2$; and there are
only $3$ jobs, of which $\{W,H\}$ is feasible, so DTLH cannot drop all $3$ (it
never drops $H$, a hopeless-free admissible job once $B$ is gone). Hence
$\mathrm{Miss}_{\mathrm{DTLH}}(G)=2$.
\end{proof}

\begin{remark}[The worst case is not forced]\label{rem:Gnotforced}
The value $2$ in Proposition~\ref{prop:gadget} \emph{does} depend on the
tie-break: processing $W$ before $B$ at $t=0$ leads DTLH to drop only one job.
Indeed, with $W$ first, $W$ is admitted ($0+3\le5$); then $B$ arrives,
$\{W,B\}$ is infeasible, and $L$ tests $s_W=3>s_B=3$ which fails, so $B$ is
dropped (miss \#1); at $t=1$ the admitted set is $\{W\}$ with remaining $2$, and
$\{W,H\}$ at $t=1$ is feasible ($1+1=2\le2$ for $H$, $2+2=4\le5$ for $W$), so $H$
is admitted with no further miss. Thus a legal DTLH execution drops only $W$'s
counterpart, matching $\mathrm{Miss}_{\mathrm{OPT}}=1$. This is precisely why we
quantify over the \emph{worst} tie-break (Remark~\ref{rem:tiebreak}): the rule as
literally stated permits the $2$-drop run, so the rule \emph{class} is at least
$2$-competitive in the worst case, even though a good tie-break is optimal on
$G$. (Verified by exhaustive enumeration of all legal tie-breaks: the drop count
on $G$ ranges over $\{1,2\}$.)
\end{remark}

\begin{theorem}[Lower bound]\label{thm:lb}
For every $k\ge1$,
\(
\mathrm{Miss}_{\mathrm{DTLH}}(\Omega_k)=2k
\)
and
\(
\mathrm{Miss}_{\mathrm{OPT}}(\Omega_k)=k,
\)
so $R\ge 2$. Moreover the bound holds for the long-run \emph{rate}: feeding the
gadget $G$ repeatedly (one copy per disjoint span of length $6$) yields a
worst-case steady miss rate of $2$ per copy under DTLH versus $1$ per copy under
OPT, a rate ratio of $2$.
\end{theorem}

\begin{proof}
The spans $[6c,6c+5]$ are pairwise disjoint and every job of $G^{(c)}$ has both
arrival and deadline inside its own span; therefore feasibility constraints
\eqref{eq:edf} never couple two different copies, and in the worst-case execution
the admitted set is empty at each epoch $6c$ (in the $2$-drop run of each copy,
$W$ is dropped and $B,H$ leave the admitted set by the end of the span).
Consequently a legal worst-case DTLH execution on $\Omega_k$ is the
concatenation of the $2$-drop executions on the individual copies, and OPT's
optimum decouples likewise. By Proposition~\ref{prop:gadget} each copy
contributes $2$ to the worst-case $\mathrm{Miss}_{\mathrm{DTLH}}$ and $1$ to
$\mathrm{Miss}_{\mathrm{OPT}}$. Summing over the $k$ copies gives $2k$ and $k$,
hence the ratio $2$. Letting $k\to\infty$ (equivalently, running the periodic
input forever) gives the stated long-run worst-case rate ratio of $2$.
\end{proof}

\section{The ratio strictly exceeds $2$: a verified instance with ratio $5$}
\label{sec:five}

The lower bound of Theorem~\ref{thm:lb} is \emph{not} tight. The following
single-busy-period instance, found by computer search and checked exhaustively,
admits a legal DTLH execution with $5$ drops while the offline optimum drops a
single job. (In the worst-case sense of Remark~\ref{rem:tiebreak}, and under the
original-size eviction convention of Remark~\ref{rem:sizeconv},
$\mathrm{Miss}_{\mathrm{DTLH}}=5$.)

\begin{proposition}\label{prop:five}
Let
\[
\omega^\star=\big\{(8,4,21),(4,5,11),(9,1,16),(8,3,14),(5,2,16),
(1,4,16),(9,1,13),(10,1,13),(7,3,18),(8,5,25)\big\}
\]
(triples $(a,s,d)$). Then $\mathrm{Miss}_{\mathrm{OPT}}(\omega^\star)=1$ and
there is a legal DTLH execution dropping $5$ jobs, so
$\mathrm{Miss}_{\mathrm{DTLH}}(\omega^\star)=5$ and
\(
R\ge 5.
\)
\end{proposition}

\begin{proof}
\emph{Offline optimum.} Deleting the single job $(4,5,11)$ leaves nine jobs, and
the following explicit single-server schedule completes all nine by their
deadlines (each interval lists the job served; the schedule is the preemptive EDF
schedule of the nine jobs, and one verifies directly that each job receives its
full size before its deadline):
\[
\begin{array}{ll}
[1,5]\ \text{serve }(1,4,16) & [11,13]\ \text{serve }(8,3,14)\\
[5,7]\ \text{serve }(5,2,16) & [13,14]\ \text{serve }(9,1,16)\\
[7,8]\ \text{serve }(7,3,18) & [14,16]\ \text{serve }(7,3,18)\\
[8,9]\ \text{serve }(8,3,14) & [16,20]\ \text{serve }(8,4,21)\\
[9,10]\ \text{serve }(9,1,13)& [20,25]\ \text{serve }(8,5,25)\\
[10,11]\ \text{serve }(10,1,13)&
\end{array}
\]
(Completion times: $(1,4,16)$ at $5\le16$, $(5,2,16)$ at $7\le16$, $(7,3,18)$ at
$16\le18$, $(8,3,14)$ at $13\le14$, $(9,1,13)$ at $10\le13$, $(10,1,13)$ at
$11\le13$, $(9,1,16)$ at $14\le16$, $(8,4,21)$ at $20\le21$, $(8,5,25)$ at
$25\le25$.) Hence $\mathrm{Miss}_{\mathrm{OPT}}(\omega^\star)\le1$. The full
ten-job set is infeasible: its preemptive EDF schedule (optimal for feasibility,
by Horn's theorem) misses a deadline, as one checks directly. Therefore
$\mathrm{Miss}_{\mathrm{OPT}}(\omega^\star)=1$.

\emph{DTLH: an explicit $5$-drop execution.} We process simultaneous $t=8$
arrivals in the order $(8,4,21),(8,3,14),(8,5,25)$ and choose victims as
indicated. At each epoch we record the current admitted set with its
\emph{remaining} sizes after EDF service up to that epoch (the server runs EDF on
the admitted set between events); the feasibility tests \eqref{eq:edf} use
remaining sizes and the eviction comparisons use original sizes
(Remark~\ref{rem:sizeconv}). Every assertion below is the displayed cumulative
EDF tableau, i.e.\ a finite list of inequalities; we write a job in the admitted
set as $(r\,@\,d)$ with $r$ its remaining size and $d$ its deadline, and order
the tableau by deadline. The cumulative entry after a row is $t$ plus the running
sum of remaining sizes up to and including that row, and feasibility means every
cumulative entry is $\le$ its deadline.

\smallskip
\emph{Admissions.} At $t=1$ admit $(1,4,16)$. After EDF service it has remaining
$1$ at $t=4$, when we admit $(4,5,11)$ (the set $\{(1@16),(5@11)\}$ is feasible
at $4$: $4+5=9\le11$, then $9+1=10\le16$). At $t=5$ admit $(5,2,16)$; at $t=7$
admit $(7,3,18)$. EDF service leaves, at $t=8$ (before the $t=8$ arrivals), the
admitted set with remaining sizes
\[
(4,5,11)\!\to\!(1@11),\quad (5,2,16)\!\to\!(2@16),\quad
(7,3,18)\!\to\!(3@18),\quad (1,4,16)\!\to\!(1@16).
\]
Now admit $(8,4,21)$ as $(4@21)$ (one checks the resulting five-job set is
feasible at $8$).

\smallskip
\emph{Miss \#1 ($t=8$, arrival $(8,3,14)$, served as $(3@14)$).} Adding it gives
the tableau at $t=8$
\[
(1@11)\!:9\le11,\ (3@14)\!:12\le14,\ (1@16)\!:13\le16,\ (2@16)\!:15\le16,\
(3@18)\!:18\le18,\ (4@21)\!:22>21,
\]
infeasible. The admitted job $(4,5,11)=(1@11)$ has original size $5>3$, and
removing it restores feasibility:
\[
(3@14)\!:11\le14,\ (1@16)\!:12\le16,\ (2@16)\!:14\le16,\ (3@18)\!:17\le18,\
(4@21)\!:21\le21.
\]
DTLH evicts $(4,5,11)$ and admits $(8,3,14)$ (miss \#1).

\smallskip
\emph{Miss \#2 ($t=8$, arrival $(8,5,25)$, $(5@25)$).} Adding it to the current
feasible set appends the row $(5@25)\!:21+5=26>25$, so the set is infeasible; and
\emph{no} admitted job has original size strictly larger than $5$, so no eviction
is allowed and $(8,5,25)$ is dropped outright (miss \#2).

\smallskip
EDF service to $t=9$ gives the admitted remaining sizes
\[
(8,3,14)\!\to\!(2@14),\ (5,2,16)\!\to\!(2@16),\ (1,4,16)\!\to\!(1@16),\
(7,3,18)\!\to\!(3@18),\ (8,4,21)\!\to\!(4@21).
\]

\emph{Miss \#3 ($t=9$, arrival $(9,1,16)$, $(1@16)$).} Tableau at $t=9$:
\[
(2@14)\!:11\le14,\ (1@16)\!:12\le16,\ (2@16)\!:14\le16,\ (1@16)\!:15\le16,\
(3@18)\!:18\le18,\ (4@21)\!:22>21,
\]
infeasible. The admitted $(1,4,16)=(1@16)$ has original size $4>1$ and its
removal restores feasibility (the running totals shift down, and the deadline-$21$
row becomes $21\le21$); DTLH evicts $(1,4,16)$ and admits $(9,1,16)$ (miss \#3).

\smallskip
\emph{Miss \#4 ($t=9$, arrival $(9,1,13)$, $(1@13)$).} Tableau:
\[
(1@13)\!:10\le13,\ (2@14)\!:12\le14,\ (2@16)\!:14\le16,\ (1@16)\!:15\le16,\
(3@18)\!:18\le18,\ (4@21)\!:22>21,
\]
infeasible. The admitted $(8,4,21)=(4@21)$ has original size $4>1$ and removing it
restores feasibility; DTLH evicts $(8,4,21)$ and admits $(9,1,13)$ (miss \#4).

\smallskip
EDF service to $t=10$ leaves the admitted remaining sizes
$(8,3,14)\!\to\!(2@14),\ (5,2,16)\!\to\!(2@16),\ (9,1,16)\!\to\!(1@16)$ and the
just-admitted $(9,1,13)$ completed by $10$ (it ran on $[9,10]$), so it leaves the
admitted set having met its deadline.

\emph{Miss \#5 ($t=10$, arrival $(10,1,13)$, $(1@13)$).} Tableau at $t=10$:
\[
(1@13)\!:11\le13,\ (2@14)\!:13\le14,\ (2@16)\!:15\le16,\ (1@16)\!:16\le16,\
(3@18)\!:19>18,
\]
infeasible. The admitted job $(8,3,14)=(2@14)$ has original size $3>1$, and its
removal restores feasibility:
\[
(1@13)\!:11\le13,\ (2@16)\!:13\le16,\ (1@16)\!:14\le16,\ (3@18)\!:17\le18 ;
\]
DTLH evicts $(8,3,14)$ and admits $(10,1,13)$ (miss \#5).

\smallskip
This is a legal DTLH run dropping the five distinct jobs
$(4,5,11),(8,5,25),(1,4,16),(8,4,21),(8,3,14)$; the surviving admitted set
$\{(5,2,16),(9,1,16),(7,3,18),(10,1,13)\}$ is feasible at $t=10$ (last tableau
above) and receives no further arrivals, so no further miss occurs. Hence
$\mathrm{Miss}_{\mathrm{DTLH}}(\omega^\star)\ge 5$. Conversely, exhaustive
enumeration of \emph{all} legal DTLH executions on $\omega^\star$ (every order of
simultaneous arrivals and every admissible victim choice, eviction by original
size) shows no execution drops more than $5$; hence
$\mathrm{Miss}_{\mathrm{DTLH}}(\omega^\star)=5$ in the worst-case sense of
Remark~\ref{rem:tiebreak}, and $R\ge 5$.
\end{proof}

\begin{remark}[The size convention matters]
Under the alternative (remaining-work) reading of the eviction test, exhaustive
simulation gives only $3$ worst-case drops on $\omega^\star$. The value $5$ uses
the original-size reading (Remark~\ref{rem:sizeconv}), which is the literal
meaning of ``size $=$ service requirement.'' A natural deterministic tie-break
(process simultaneous arrivals in increasing size, always evict the largest
eligible job) drops fewer than $5$ on $\omega^\star$; the value $5$ is realized
by the adversarial tie-break above, which is why we measure the worst case.
\end{remark}

\begin{remark}[Mechanism]
The gap is created by DTLH's \emph{myopic swap rule}. When a small job arrives
into an infeasible queue, DTLH may sacrifice a strictly larger admitted job---a
job which may already have absorbed substantial server time, so its eviction also
\emph{wastes} that work. Across a burst of small arrivals this can discard
several distinct large jobs, one per arrival, whereas the clairvoyant optimum
removes a \emph{single} well-chosen ``linchpin'' job (typically the
earliest-deadline one) whose deletion restores feasibility for everything else
and never wastes service on a job it will drop. Because the wasteful swaps can be
concentrated in one busy period while the offline fix is a single deletion,
$\mathrm{Miss}_{\mathrm{DTLH}}$ can exceed $\mathrm{Miss}_{\mathrm{OPT}}$ by a
multiplicative factor strictly larger than $2$.
\end{remark}

\section{Status of the exact value of $R$}

Propositions \ref{prop:gadget}, \ref{prop:five} and Theorem~\ref{thm:lb} prove
\[
R\ \ge\ 5 .
\]
Beyond this, the exact value of $R$ is \emph{not} resolved here. We summarize
honestly what the evidence shows and where it stops.

\begin{remark}[Verified evidence that $R$ keeps growing]\label{rem:grow}
For each target value $v\in\{2,3,4,5\}$ we exhibit a single-busy-period instance
$\omega_v$ with $\mathrm{Miss}_{\mathrm{OPT}}(\omega_v)=1$ and worst-case
$\mathrm{Miss}_{\mathrm{DTLH}}(\omega_v)=v$, where $\mathrm{Miss}_{\mathrm{OPT}}$
is computed exactly by subset enumeration and $\mathrm{Miss}_{\mathrm{DTLH}}$ by
exhaustive enumeration of all legal DTLH tie-breaks (arrival orders and victim
choices, original-size eviction). Concretely:
\begin{align*}
\omega_2&=\{(5,1,18),(5,5,16),(6,4,14),(7,2,10),(11,2,15),(11,1,13),(8,5,28),(9,4,23),(8,2,25)\},\\
\omega_3&=\{(8,5,26),(10,2,27),(11,3,22),(11,1,24),(8,5,26),(2,3,10),(7,4,15),(8,1,21),(4,4,23),(11,1,21)\},\\
\omega_4&=\{(1,3,18),(8,5,26),(0,5,16),(6,5,11),(1,4,6),(11,1,18),(10,2,18),(12,2,23),(12,2,18),(9,2,25)\},\\
\omega_5&=\omega^\star\text{ of Proposition~\ref{prop:five}}.
\end{align*}
All four equalities $\mathrm{Miss}_{\mathrm{OPT}}=1$ and the corresponding
worst-case $\mathrm{Miss}_{\mathrm{DTLH}}=v$ were machine-verified. Thus a
\emph{single} value of $\mathrm{Miss}_{\mathrm{OPT}}$ already supports
arbitrarily large \emph{verified} numerators up to $5$, and a local search over
larger instances with $\mathrm{Miss}_{\mathrm{OPT}}=1$ continues to increase the
attainable worst-case DTLH miss count without revealing any constant ceiling.
This strongly suggests
\[
R=\sup_\omega
\frac{\mathrm{Miss}_{\mathrm{DTLH}}(\omega)}{\mathrm{Miss}_{\mathrm{OPT}}(\omega)}
=\infty,
\]
i.e.\ that the DTLH policy class is \emph{not} constant-competitive against the
clairvoyant offline optimum, with the obstruction being the
one-larger-job-per-arrival swap rule (which both discards and wastes work).
However, we have \emph{not} produced a closed-form infinite family that provably
forces
$\mathrm{Miss}_{\mathrm{DTLH}}/\mathrm{Miss}_{\mathrm{OPT}}\to\infty$, and we do
\emph{not} claim $R=\infty$ as a theorem.
\end{remark}

\begin{remark}[Honest summary of what is proved]
\begin{enumerate}
\item (Lemma~\ref{lem:nowaste}) DTLH is optimal whenever the input is
schedulable ($\mathrm{Miss}_{\mathrm{OPT}}=0\Rightarrow
\mathrm{Miss}_{\mathrm{DTLH}}=0$, for every tie-break): it never drops a job
needlessly.
\item (Lemma~\ref{lem:onewindow}) On any single common release/deadline window
(the ``keep-the-most-jobs'' knapsack-by-count regime), the \emph{evict-largest}
DTLH rule $\mathrm{DTLH}_{\max}$ is exactly optimal. This is \emph{not} true for
every tie-break: the worst-case tie-break already loses a factor (a single window
with capacity $5$ and sizes $3,2,1,2$ gives
$\mathrm{Miss}_{\mathrm{DTLH}}=2>1=\mathrm{Miss}_{\mathrm{OPT}}$,
Remark~\ref{rem:onewindowbad}), so DTLH's competitive loss stems from the swap
rule's victim choice and is visible even in one window.
\item (Theorem~\ref{thm:lb}) $R\ge 2$, via an explicit scalable family that also
yields a long-run worst-case miss-rate ratio of $2$. (The bound is for the
worst-case tie-break; a good tie-break is optimal on the gadget, see
Remark~\ref{rem:Gnotforced}.)
\item (Proposition~\ref{prop:five}) $R\ge 5$, via a fully verified instance,
exhibited together with the explicit legal DTLH execution that realizes $5$
drops (worst-case tie-break, original-size eviction).
\end{enumerate}
The remaining open point is the exact value of $R$: the verified evidence
(Remark~\ref{rem:grow}) points to $R=\infty$ (no constant competitive ratio),
with the obstruction being DTLH's one-larger-job-per-arrival swap rule, but a
closed-form unbounded construction is left unresolved. We also note that all of
the above is for the \emph{worst-case} reading of the underspecified policy
(Remark~\ref{rem:tiebreak}); the competitive ratio of any \emph{fixed}
deterministic tie-break could be strictly smaller and is likewise open.
\end{remark}

\end{document}
